\section{Execution Flow}

The championship progresses sequentially from round to round, while the matches inside each round run in parallel.

\begin{enumerate}
    \item \texttt{PlayChampionship} starts the tournament and keeps the current player list.
    \item \texttt{PlayRound} splits that list into pairs and launches one \textbf{goroutine} per \textbf{match}.
    \item Each worker calls \texttt{PlayMatch} exactly once.
    \item The round coordinator waits for one outcome per match.
    \item Results are stored by \texttt{matchIndex}, so the bracket order remains stable even if workers finish out of order.
    \item The winners returned by the round become the input of the next round.
\end{enumerate}

The worker logic is intentionally small.
It calls the toss function exactly once, maps \texttt{Heads} to the first player and \texttt{Tails} to the second player, and returns a
\texttt{MatchResult} with the round number, match number, participants, winner, and tossed side.

The round coordinator is the \textbf{synchronization point}.
It does not poll or guess completion; it waits for exactly one result per started match.
That is the key mechanism that keeps the implementation concurrent inside each round and deterministic at the bracket level.

\medskip

\noindent\textbf{Concurrent match launch in \texttt{src/championship/round/round.go}}
\begin{lstlisting}[style=reportcode]
for matchIndex := 0; matchIndex < matchCount; matchIndex++ {
    matchNumber := matchIndex + 1
    firstPlayer := players[matchIndex*2]
    secondPlayer := players[matchIndex*2+1]

    go func(matchIndex, matchNumber int, firstPlayer, secondPlayer domain.Player) {
        result, err := match.PlayMatch(roundNumber, matchNumber, toss, firstPlayer, secondPlayer)
        outcomes <- matchOutcome{matchIndex: matchIndex, result: result, err: err}
    }(matchIndex, matchNumber, firstPlayer, secondPlayer)
}
\end{lstlisting}

\noindent\textbf{Round orchestration in \texttt{src/championship/tournament/championship.go}}
\begin{lstlisting}[style=reportcode]
for len(currentPlayers) > 1 {
    winners, matches, err := round.PlayRound(roundNumber, currentPlayers, toss)
    if err != nil {
        return domain.ChampionshipResult{}, err
    }

    roundResults = append(roundResults, domain.NewRoundResult(roundNumber, matches, winners))
    currentPlayers = winners
    roundNumber++
}
\end{lstlisting}
